\section{Introduction and statements}\label{sec:introduction}

The Kuga--Satake construction associates an abelian variety to a polarized Hodge structure of K3 type. Its original purpose was to place the second cohomology of a K3 surface inside a tensor construction on the first cohomology of an abelian variety \cite{KugaSatake}. The resulting embedding is a Hodge morphism. Its realization by an algebraic correspondence is a further assertion, usually called the algebraicity of the Kuga--Satake correspondence; see \cite[Chapter 4]{Huybrechts} and \cite{vanGeemen}. This assertion is a particular instance of the rational Hodge conjecture.

We prove two conditional reductions for this algebraicity problem. At a point where the primitive cohomology has full orthogonal Hodge group, an algebraic correspondence inducing any nonzero map from primitive cohomology to the second cohomology of an abelian variety suffices to recover the \emph{specified} map into the \emph{full} even-Clifford target. Very-general existence of such a map then implies algebraicity of the specified correspondence at every point of the same polarized component. These reductions do not establish the existence of the initial map and therefore do not resolve the unrestricted Kuga--Satake conjecture.

The proof uses three devices which may be useful in related algebraicity questions. Independent operations on the two weight-one factors recover every multiplicity, and rational descent handles nonsplit Clifford algebras. Proper Hilbert parameter spaces specialize the \emph{entire prescribed flat class}, including its rational coefficients, over an analytic period chart. Finally, signed square roots compare quadratic polarizations through operations on the abelian factors. This last step does not require algebraicity of a Hodge endomorphism of the K3 surface.

\subsection{Conventions and the exact target}

All varieties in the statements are over $\mathbb C$, and cohomology is rational Betti cohomology unless a different field is displayed. We write $\operatorname{CH}^r(-)_{\mathbb Q}$ for codimension-$r$ cycles modulo rational equivalence with rational coefficients. A subscript records dimension, so $\operatorname{CH}_a$ denotes cycles of dimension $a$. We write $\operatorname{NS}(S)$ for the N\'eron--Severi group, $\operatorname{Hg}$ for the Hodge group, $\operatorname{Hom}_{\mathrm{HS}}$ for rational Hodge maps, and $F^\bullet$ for the Hodge filtration. The Tate twist $(1)$ shifts Hodge bidegrees by $(-1,-1)$. On a Hodge component of type $(p,q)$ the Hodge circle acts as $e^{i(p-q)\theta}$. The joint Hodge group of several structures is the smallest rational algebraic group containing their simultaneous Hodge circle. Equal-weight rational maps are Hodge maps exactly when they intertwine this circle.

Let $(S,h)$ be a primitively amply polarized K3 surface, with $h^2=2d$. Put
\[
 V=h^\perp\subset H^2(S,\mathbb Q),\qquad
 T=T(S)=\operatorname{NS}(S)_{\mathbb Q}^{\perp}\subset V.
\]
Orthogonality here is for the cup form $q$. The space $V$ has dimension $21$ and signature $(2,19)$, with the positive two-plane determined by $H^{2,0}(S)$. The space $T$ has the same period line and dimension at least two.

For the final target, fix a rational quadratic polarization $Q$ of $T$ and all standard Kuga--Satake data: the full even Clifford construction, its parity identification by right multiplication, a rational weight-one polarization, and any chosen isogenous realization. Denote this abelian variety by $A_S$, and set $W_S=H^1(A_S,\mathbb Q)$. The standard map $T(1)\to\operatorname{End}(W_S)$ and the stipulated identification $W_S^\vee\simeq W_S(1)$ determine
\begin{equation}\label{eq:target}
 \kappa_S:T\longrightarrow W_S\otimes W_S
 \lhook\joinrel\longrightarrow H^2(A_S\times A_S,\mathbb Q).
\end{equation}
The last arrow is the ordered K\"unneth inclusion. To realize this map means to find
\begin{equation}\label{eq:cycle-target}
 \Gamma_S\in\operatorname{CH}^2(S\times A_S\times A_S)_{\mathbb Q},
 \qquad
 (p_{AA})_*\bigl(\operatorname{cl}(\Gamma_S)\cup p_S^*x\bigr)
 =\kappa_S(x)\quad(x\in T).
\end{equation}
Here $\operatorname{cl}$ is the rational Betti cycle class, and $p_S$ and $p_{AA}$ are the projections to $S$ and $A_S\times A_S$. Every normalization in \eqref{eq:target} is retained. The target is never replaced by a simple spin factor.

\subsection{The conditional theorems}

\begin{theorem}[Recovery at a point with full Hodge group]\label{main:generic}
Suppose $\operatorname{Hg}(V)=\operatorname{SO}(V,q)$. The following conditions are equivalent.
\begin{enumerate}
\item There are an abelian variety $B$ and a cycle $\alpha\in\operatorname{CH}^2(S\times B)_{\mathbb Q}$ such that
\[
 \alpha_*:V\longrightarrow H^2(B,\mathbb Q)
\]
is nonzero.
\item There are a smooth connected projective surface $Y$ and a dominant generically finite morphism $f:Y\to S$ such that
\begin{equation}\label{eq:wedge-condition}
 f_*\!\left(\operatorname{im}\!\left[
 \bigwedge\nolimits^2 H^0(Y,\Omega_Y^1)
 \longrightarrow H^0(Y,\Omega_Y^2)\right]\right)\ne0.
\end{equation}
\item The prescribed full Kuga--Satake map is realized as in \eqref{eq:cycle-target}, for any stipulated standard data on $T$.
\end{enumerate}
In this situation $T=V$. The same input realizes the standard map for the full primitive Clifford construction with the cup form.
\end{theorem}

In \eqref{eq:wedge-condition}, $f_*$ is the trace on holomorphic two-forms, equivalently the restriction of cohomological pushforward to $H^{2,0}$. There is no genericity condition on $B$. The hypothesis on $V$ holds outside a countable union of proper closed analytic subsets in every full polarized period chart; see Proposition~\ref{gen:genericity}.

\begin{theorem}[Specialization in a polarized component]\label{main:specialization}
Fix a component of the moduli space of primitively amply polarized K3 surfaces of degree $2d$. Assume that condition \textup{(i)} or condition \textup{(ii)} of Theorem~\ref{main:generic} holds very generally in that component. Then every $(S,h)$ in the component satisfies \eqref{eq:cycle-target}, with all the stipulated standard data on $T(S)$.

Here it suffices that, about each desired point, the input holds outside a countable union of proper closed analytic subsets of a full polarized period chart. In particular an input outside countably many proper algebraic subsets of the moduli component suffices. The witness cycles, abelian varieties, or surface covers need not vary in a family.
\end{theorem}

\begin{corollary}\label{main:all-degrees}
If the hypothesis of Theorem~\ref{main:specialization} holds in the required polarized components in every degree, then \eqref{eq:cycle-target} holds for every projective K3 surface.
\end{corollary}

\begin{proof}
Every projective K3 surface has an ample integral class. Divide it to primitivity in $H^2(S,\mathbb Z)$; the resulting integral $(1,1)$ class is the first Chern class of a line bundle, a positive power of which is ample. Hence it is itself ample, and its square is a positive even integer. Apply Theorem~\ref{main:specialization} in its component.
\end{proof}

An input in one polarized component yields the conclusion in that component. Constructing a very-general correspondence or surface cover with the stated property remains the external existence problem. The hypotheses concern actual algebraic cycles; absolute Hodge or motivated classes do not supply this input.

\subsection{Proof outline}

Section~\ref{sec:cycles} proves the equivalence of the two inputs and records the available cycle operations. Section~\ref{cliff:section} fixes the full Clifford target and tensor conventions. Section~\ref{gen:section} recovers that target from one nonzero input by identifying the spin constituents and applying independent multiplicity operations. Section~\ref{spec:section} specializes the entire prescribed class in a local projective family. Section~\ref{data:section} restricts to $T$ and compares quadratic polarizations; Section~\ref{sec:conclusion} assembles the two theorems. Appendix~\ref{app:dependencies} records the established inputs and their roles.
